% José Gabriel Maza Mendoza — CV (ES)
% Plantilla inspirada en layout LaTeX de referencia (Tony Gualdrón)
\documentclass[10.5pt,a4paper]{article}

\usepackage[margin=1.55cm]{geometry}
\usepackage{fontspec}
\usepackage{xcolor}
\usepackage{hyperref}
\usepackage{enumitem}
\usepackage{titlesec}
\usepackage{setspace}
\usepackage{tabularx}
\usepackage{fontawesome5}

\setmainfont{Times New Roman}
\setsansfont{Arial}

\definecolor{ink}{HTML}{111111}
\definecolor{muted}{HTML}{333333}
\definecolor{sectionline}{HTML}{222222}

\hypersetup{
  colorlinks=true,
  urlcolor=ink,
  pdfauthor={José Gabriel Maza Mendoza},
  pdftitle={José Gabriel Maza Mendoza — CV}
}

\pagestyle{empty}
\setlength{\parindent}{0pt}
\setlength{\parskip}{0.25em}
\setlist[itemize]{leftmargin=1.15em, itemsep=0.12em, topsep=0.2em, parsep=0pt}
\setstretch{1.08}

\titleformat{\section}
  {\normalsize\bfseries\sffamily}
  {}{0pt}{\MakeUppercase}[\vspace{-0.35em}{\color{sectionline}\titlerule[0.7pt]}\vspace{0.35em}]
\titlespacing*{\section}{0pt}{0.85em}{0.35em}

\newcommand{\Job}[4]{%
  \noindent\textbf{#1}\hfill {\small #3}\\[-0.05em]
  \textit{#2}\hfill {\small\color{muted}#4}\\[0.15em]
}

\newcommand{\Edu}[4]{%
  \noindent\textbf{#1}\hfill {\small #3}\\[-0.05em]
  \textit{#2}\hfill {\small\color{muted}#4}\\[0.25em]
}

\newcommand{\Award}[4]{%
  \noindent\textbf{#1}\hfill {\small #3}\\[-0.05em]
  \textit{#2}\\[0.05em]
  {\small #4}\\[0.35em]
}

\begin{document}

\begin{center}
  {\LARGE\bfseries José Gabriel Maza Mendoza}\\[0.25em]
  {\large\textit{Ingeniero de Sistemas · Fullstack / Product Engineer}}\\[0.45em]
  {\small
    \faEnvelope\ \href{mailto:mazamendoza123@gmail.com}{mazamendoza123@gmail.com}
    \quad \faPhone\ (+57) 300 8420902
    \quad \faMapMarker*\ Barranquilla, Colombia
  }\\[0.15em]
  {\small
    \faGithub\ \href{https://github.com/Jgmaza}{github.com/Jgmaza}
    \quad \faGlobe\ \href{https://jgmaza.vercel.app}{jgmaza.vercel.app}
  }
\end{center}

\section{Perfil profesional}
Ingeniero de Sistemas y desarrollador fullstack con enfoque en product engineering.
Conecto interfaces, APIs, datos e infraestructura para llevar soluciones desde la
definición del problema hasta el despliegue. Experiencia en desarrollo web,
integraciones entre sistemas, herramientas internas y asistentes de IA, usando
React, Next.js, TypeScript, NestJS, FastAPI y Supabase según las necesidades de
cada proyecto. Cómodo con el ciclo completo: diseño, implementación, calidad y deploy.

\section{Experiencia laboral}

\Job{Desarrollador Fullstack}{Ministerio de Minas y Energía}{Sep 2025 -- presente}{Barranquilla, Colombia}
Responsabilidades: desarrollo y mantenimiento en paralelo de varias líneas productivas
institucionales --- integrador de servicios / gateway de interoperabilidad, plataforma
de vigilancia energética (dosimetría) y aplicación de mapeo de stakeholders --- con foco
en APIs confiables, consistencia de datos y frontends usables bajo restricciones
institucionales.
\begin{itemize}
  \item Expandí las interoperabilidades del integrador de \textbf{2 a 6} rutas/familias
        de integración sobre un gateway HTTP (APISIX + FastAPI + Vue).
  \item Unifiqué dosimetría end-to-end (API Node, SQL Server, Angular) sobre volúmenes
        de prueba del orden de \textbf{$\sim$629k} registros históricos y
        \textbf{$\sim$61k} nuevos; CLI/scripts de repair recuperaron y sincronizaron
        datos históricos hacia la capa unificada.
  \item Aceleré el cierre de entregas frente al ritmo previo reportado por el área
        (el trabajo venía retrasado con el desarrollador anterior), sosteniendo varias
        HU en paralelo con despliegues documentados a QA.
\end{itemize}

\vspace{0.35em}
\Job{Desarrollador Fullstack}{Guarapo Labs}{Jun 2025 -- Ago 2025}{Barranquilla, Colombia}
Apoyé entregas a cliente internacional y un producto nuevo (front, back e IA).
Construí herramientas internas con Next.js y Supabase.
\begin{itemize}
  \item Al incorporar me, aceleramos entregas y nos pusimos al día con el cronograma;
        en el frente visual mantuve las tareas de UI al corriente.
\end{itemize}

\vspace{0.35em}
\Job{Desarrollador Fullstack · IA / automatización}{Clazz Dev}{Feb 2025 -- May 2025}{Barranquilla, Colombia}
Diseñé automatizaciones e asistentes con OpenAI, n8n, FastAPI y Python --- agentes
dentro de flujos reales (ventas y operaciones), no solo chat.
\begin{itemize}
  \item Workflows de agente de ventas que redujeron la operación humana de aproximadamente
        \textbf{7 a 2 personas}, al gestionar el ciclo de ventas de punta a punta.
\end{itemize}

\vspace{0.35em}
\Job{Desarrollador Fullstack}{I-Cluster}{Jun 2024 -- Ene 2025}{Barranquilla, Colombia}
Desarrollé software para una iniciativa municipal de salud emocional (Tolú) con React
y Express. Colaboración transversal en features orientadas a programas de bienestar,
con énfasis en fiabilidad y entrega mantenible.

\vspace{0.35em}
\Job{Desarrollador Fullstack}{Infinity Tech Consulting}{Abr 2022 -- Jun 2024}{Barranquilla, Colombia}
Construí y mantuve aplicaciones web fullstack con Next.js, Flask y Prisma --- frontend
y backend --- desplegadas en AWS y Vercel.
\begin{itemize}
  \item \textbf{Commodo}: sistema de fidelización (estilo puntos) que sigue en producción;
        administré el front de las \textbf{3 plataformas web} del cliente con cambios
        del día a día.
  \item POS desde cero en $\approx$\textbf{1 mes}, integrando APIs propias de Commodo,
        Siigo y consultas externas (p.\ ej.\ departamentos).
\end{itemize}

\vspace{0.35em}
\Job{Soporte / Freelance (MIOBOX)}{Sofin SAS}{Nov 2021 -- Ene 2022; freelance $\approx$2 meses en 2026}{Barranquilla, Colombia}
Capítulo inicial: soporte de facturación electrónica y pruebas de código ERP.
Retorno freelance: puente NestJS SOAP/REST entre ERP y operaciones de planta
(órdenes, consumo, cierre de entrega).
\begin{itemize}
  \item Reduje el ciclo de actualización en el ERP de aproximadamente
        \textbf{40 minutos a $\sim$2 minutos} mediante el puente de integración.
\end{itemize}

\section{Educación formal}
\Edu{Ingeniero de Sistemas}{Universidad del Norte}{2018 -- 2022}{Barranquilla, Colombia}
Distinción: Estudiante Destacado.

\section{Habilidades}
\textbf{Core:} TypeScript, React, JavaScript, Java, Python, NestJS, FastAPI, Node/Express,
REST APIs, SQL (MySQL, PostgreSQL, SQL Server), Git.\\
\textbf{También en producción:} Next.js, Vue 3, Angular, Prisma, Supabase, AWS, Vercel,
APISIX / API gateways, integraciones SOAP, n8n, HTML/CSS/Tailwind, Linux.

\section{Reconocimientos y premios}
\Award{1.\textsuperscript{er} puesto --- Interledger Hackathon}{Ledger One}{2026}{%
  Primer puesto con Ledger One, proyecto de equipo sobre open payments / Interledger.}
\Award{Mejor proyecto con ShadAI · 4.\textsuperscript{o} lugar general}{BarranquillIA 2025 · EduNova}{2025}{%
  Premio a mejor proyecto usando ShadAI (tecnología USA) y 4.\textsuperscript{o} lugar
  general. Plataforma educativa en reconstrucción de producto.}

\section{Idiomas}
Inglés --- EF SET B2 \quad|\quad Español --- Nativo

\section{Certificaciones y cursos}
\begin{itemize}
  \item \textbf{Blockchain Builders 101} --- Ruta N / Stellar Development Foundation / BAF
        (en curso; bootcamp 5 semanas, MVP sobre Stellar).
  \item Certificate of Participation --- ``Build with AI'', Google Developers Group
        Barranquilla · 2025.
  \item React Query --- DevTalles · 12/2023; HTML5 y CSS3 --- Udemy · 07/2018.
\end{itemize}

\vspace{0.6em}
{\footnotesize\color{muted}Última actualización: octubre 2026}

\end{document}
